\preliminar{Resumen del proyecto}

\textit{Diffusion Policy} genera acciones de manipulaci\'on rob\'otica mediante un proceso
de difusi\'on condicionado por la observaci\'on. Su variante visuomotora encadena un
codificador visual y una red generativa, de modo que la representaci\'on de la imagen
delimita la informaci\'on disponible para el control. La investigaci\'on posterior ha
modificado sobre todo el bloque generativo, mientras que la elecci\'on del codificador y de
su estrategia de entrenamiento sigue sin resolverse cuando solo se dispone de un conjunto
reducido de demostraciones.

Este trabajo eval\'ua la influencia del codificador visual y de su estrategia de
entrenamiento sobre el rendimiento de una \textit{Diffusion Policy} en la tarea Push-T. Con
ese fin se implementan cinco variantes del codificador dentro de la misma pol\'itica: una
ResNet-18 entrenada desde cero, la misma red preentrenada sobre ImageNet en versiones
congelada y con ajuste fino, y dos transformadores de visi\'on congelados, DINOv2 ViT-S/14
y CLIP ViT-B/16. Todas entregan un vector de \num{512} componentes y comparten red
generativa, hiperpar\'ametros, datos y protocolo de evaluaci\'on, de manera que la \'unica
variable manipulada es el codificador. El ajuste emplea \num{206} demostraciones
teleoperadas y un presupuesto m\'aximo de \num{500} \'epocas con parada anticipada. La
evaluaci\'on mide la puntuaci\'on media de cobertura sobre \num{50} episodios generados
con semillas reservadas, junto con el tiempo de c\'omputo consumido.

Los resultados corresponden a las tres variantes convolucionales completadas, ya que las
dos basadas en transformador siguen en ejecuci\'on. La l\'inea base entrenada desde cero
alcanza una puntuaci\'on de \num{0.864}, frente a \num{0.668} de la variante congelada y
\num{0.648} de la ajustada, ca\'idas del \SI{23}{\percent} y del \SI{25}{\percent} cuyos
intervalos de confianza excluyen el cero. El
preentrenamiento tampoco compensa en coste: cada \'epoca de la l\'inea base ocupa
\SI{1.5}{\minute}, frente a \SI{5.3}{\minute} y \SI{14.7}{\minute} de las variantes
preentrenadas, y las tres ejecuciones suman casi \num{200} horas de procesamiento
gr\'afico. Ese sobrecoste se atribuye a la resoluci\'on de entrada y no al n\'umero de
par\'ametros entrenables, que apenas var\'ia entre variantes.

Se concluye que, en Push-T y con pocas demostraciones, el preentrenamiento supervisado
sobre im\'agenes naturales degrada el rendimiento y encarece el entrenamiento. La ventaja
de la l\'inea base debe atribuirse al codificador junto con su cadena de preprocesado,
puesto que ambos var\'ian de forma conjunta.

\vspace{1em}
\noindent\textbf{Palabras clave:} aprendizaje por imitaci\'on, modelos de difusi\'on,
codificador visual, transferencia de aprendizaje, manipulaci\'on rob\'otica.
